\section{Conclusion}
\label{sec:conclusion}

A unified drift-diffusion--Poisson model of the near-electrode layers of an atmospheric-pressure argon arc has been presented. A single set of equations is solved from the electrode surface, through the space-charge sheath and the ionization layer, to the bulk plasma: the continuity equations with fluxes derived from the species momentum balances, including thermal diffusion and the collisional drag between species, separate electron and heavy-particle energy equations, a heavy-gas momentum equation and Poisson's equation. Adaptive mesh refinement resolves the sheath on a fraction of the Debye length. The surface temperature of the cathode follows from a heat balance that includes thermionic emission with Schottky lowering. The model reproduces the near-cathode layer computed by Almeida et al.\ \cite{Almeida2008} and the near-anode layer computed by Semenov et al.\ \cite{semenov2016numerical}: the temperatures agree within 3--5\,\%, the charged-particle densities within 15\,\% outside the cathode sheath, and the ion fraction of the current at the surface within 5\,\%.

For a cathode at \(T_{\rm bulk}=21\,000\)~K and \(10^{8}\ \mathrm{A\,m^{-2}}\), and an anode at \(15\,000\)~K and \(4\times10^{6}\ \mathrm{A\,m^{-2}}\), with the radial heat loss of the column evaluated for assumed column radii of 0.5~mm (cathode) and 3.5~mm (anode), the main results are as follows.
\begin{itemize}
\item The two layers have opposite structures. The cathode has an ion-rich sheath \(0.58\ \mu\)m wide with a voltage of 3.6~V, followed by an ionization layer of about \(10^{-4}\)~m in which the gas is heated, the plasma becomes almost fully ionized and the charged-particle density, driven by the hot electrons, approaches its equilibrium value at the heavy-particle temperature only at the end of the layer. Joule heating is the dominant electron energy source throughout. The anode has a thin ion-rich sheath (1.4~V across \(1.0\ \mu\)m). Most of its 5.2~V retarding drop is spread over the quasi-neutral layer, which extends over the whole domain as the plasma cools from the imposed bulk state; the charged particles are supplied by diffusion at densities far above ionization equilibrium, and \(T_e-T_h\) reaches 5.1~kK at the surface. The electrons lose energy to the retarding field, and this loss, together with the radiation and the radial loss of the bulk plasma, is supplied by the electron energy flux from the column.
\item The electrons cross both sheaths with fewer than 0.25 collisions, yet their density follows the Boltzmann relation within about 2\,\%, so that the electron force balance remains adequate there. The ions cross 1.8 (cathode) and 9.9 (anode) mean free paths in the sheaths and do not satisfy the Bohm criterion at the sheath edge, where they move at 0.78 and 0.19 times the Bohm velocity. A Bohm flux at the sheath edge would overestimate the ion flux by 28\,\% at the cathode and by a factor of 5 at the anode.
\item The cathode heat flux, \(27\ \mathrm{MW\,m^{-2}}\), is the small difference between the heating of the surface (\(566\ \mathrm{MW\,m^{-2}}\)) and emission cooling (\(539\ \mathrm{MW\,m^{-2}}\)). Because emission cooling depends steeply on \(T_{\rm wall}\), the heat flux into the cathode body is set by the electrode rather than by the plasma: a 5\,\% change in the heating of the surface would shift \(T_{\rm wall}\) by only about 20~K and the heat flux by less than 1\,\%. The anode heat flux, \(36\ \mathrm{MW\,m^{-2}}\), is 66\,\% electronic and 31\,\% gas conduction. Over \(2\times10^{6}\)--\(10^{7}\ \mathrm{A\,m^{-2}}\) it increases linearly with the current: the electron terms add 5.6--6.3~W per ampere, of which 4.7--4.9~W per ampere is the work function (\(W=4.65\)~eV, copper), while gas conduction, \(11\)--\(12\ \mathrm{MW\,m^{-2}}\), is set by the bulk temperature and not by the current.
\item Because the ions cross the cathode sheath with only one to two collisions, a local estimate of their impact energy is not reliable, and only its maximum, for an ion crossing the sheath without a collision, has been given: 4.9~eV, corresponding to at most \(26\ \mathrm{MW\,m^{-2}}\). Owing to the stiffness of emission cooling, this term would change the cathode surface temperature by only about 20~K.
\item Since the anode attachment can be constricted or diffuse, the anode was studied over a range of bulk temperatures and current densities. The ion current density into the surface hardly depends on the arc current, so that the ion fraction of the current falls from 6.4 to 1.0\,\% between \(2\times10^{6}\) and \(10^{7}\ \mathrm{A\,m^{-2}}\); the anode sheath width decreases from 1.16 to \(0.77\ \mu\)m over this range, and from 1.13 to \(0.91\ \mu\)m as the bulk temperature increases from 13\,000 to 17\,000~K.
\item After a change of the current, the cathode layer reaches a new steady state within about \(3\ \mu\)s, relaxing with \(\tau=0.36\ \mu\)s. The anode layer relaxes with \(\tau=8\)--\(9\ \mu\)s over all the anode cases, almost independently of the current density; most of the change of its heat flux, which is carried by the electron current, occurs within 0.1~ns, and both the heat flux and the ion current settle within about \(30\ \mu\)s.
\end{itemize}

The local, inertialess transport description is marginal within about \(0.3\ \mu\)m of the cathode, where the inertia of the ions exceeds their drag. In the same region, the field heats the ions well above the atom temperature, so that the single heavy-particle temperature does not describe their random energy, and the ion energy distribution at the cathode requires a kinetic treatment of the sheath. Elsewhere, ions and atoms exchange energy over distances short compared with the gradients of \(T_h\), and a single heavy-particle temperature is adequate. The present one-dimensional calculations represent the radial losses of the column through a sink for an assumed column radius. Since the layers settle within microseconds at the cathode and tens of microseconds at the anode, and are characterized by the current density and the state of the adjacent column, the model can supply electrode boundary conditions, such as the voltage drop, the current partition and the heat flux to the electrode, to LTE models of the arc column.
